\documentclass[14pt, a4paper]{article}

\usepackage{fontspec}
\usepackage{polyglossia}
\usepackage{etoolbox}
\setmainlanguage{russian}
\setotherlanguage{english}

\setmainfont{Times New Roman}
\newfontfamily\cyrillicfont{Times New Roman}
\newfontfamily\cyrillicfontsf{Times New Roman}
\newfontfamily\cyrillicfonttt{Times New Roman}

\usepackage[bottom]{footmisc}
\usepackage{amsmath}
\usepackage{enumitem}
\usepackage{geometry}
\usepackage{booktabs}
\usepackage{array}
\usepackage{float}
\usepackage{graphicx}
\usepackage{multicol}
\usepackage{hyperref}
\usepackage{bigfoot}
\DeclareNewFootnote{default}


\geometry{margin=0.5in}
\makeatletter
\renewcommand{\normalsize}{\fontsize{12pt}{14.4pt}\selectfont}
\normalsize
\makeatother

\setlength{\columnsep}{0.9cm}
\setlength{\columnseprule}{0pt}
\setlength{\parindent}{3em}

\title{\textbf{Фермеры как социально-профессиональная группа}\\[0.35em]
	\large Декомпозиция Блайндера–Оаксаки социально-демографических признаков на основе теории Бурдье}
\author{Матвеев Д.\,А.}
\date{19 октября 2026 г.}

\begin{document}
	\renewcommand{\footnoterule}{}
	\maketitle
	
	%%\begin{abstract}
	%%	bla bla bla
	%%\end{abstract}
	
	%%\vspace{1em}
	
	\section{Введение}\label{sec:intro}
	Фермеры как социальная группа преодолели значительную эволюцию со времен распада СССР. Как отмечает Боконтаева Д. К. \cite{bokontaeva1998farmers} основным драйвером переходного процесса с колхозного типа хозяйства к фермерскому послужил переход к рыночной экономике. В частности, исследователь отмечает роль предпринимательства как основа перехода профессиональной группы. Однако стремление к предпринимательству влечет за собой множество издержек, которые значительно влияют на профессиональные траектории фермеров. 
	
	Как отмечает Великий П. П.  \cite{velikiy2024fermerstvo}, российские фермеры не являются до конца «укорененным» социальным классом. Они действуют в условиях высокой экономической неопределенности, демонстрируя модели поведения и ведения бизнеса, отличные от тех, что свойственны остальному сельскому населению. Их часто воспринимают как наиболее благополучный сегмент сельской рабочей силы, однако их экономическое положение остается структурно хрупким.
	
	Для понимания экономического положения фермеров предлагаемое исследование опирается на теорию капитала и поля Пьера Бурдье. Бурдье \cite{bourdieu1986forms} утверждает, что капитал существует в трех основных формах: экономический капитал (финансовые ресурсы), культурный капитал (образование, дипломы, навыки) и социальный капитал (сети отношений). Ключевой тезис Бурдье заключается в том, что ценность капитала не является внутренне присущей; она определяется специфической социальной ареной, на которой действуют акторы: «поле». Один и тот же культурный капитал (например, университетский диплом) может высоко вознаграждаться в одном поле (например, в городской корпоративной занятости), но обесцениваться или не признаваться в другом (например, в сельскохозяйственном самозанятости).
	
	Центральная проблема крестьяноведения и экономики труда заключается в понимании того, является ли экономическое благополучие фермеров результатом композиционных эффектов (различий в наблюдаемых характеристиках, таких как возраст и образование) или структурных эффектов (различий в отдаче от этих характеристик, отражающих «\textit{несоответствие поля}», описанное Бурдье).
	
	Если разрыв в доходах в первую очередь обусловлен дефицитом человеческого капитала, меры социальной политики должны быть сосредоточены на сельском образовании и профессиональной подготовке. Однако если разрыв вызван несоответствии полей по Бурдье, то проблема кроется в структурной организации сельскохозяйственного рынка труда. Исследование направлено на разграничение этих двух механизмов.
	
	Для определения природы такого разрыва, в исследовании предлагается оценка различий в:
	\begin{enumerate}[label=\arabic*)]
		\item \textit{социально-демографическом профиле} фермеров и не-фермеров в Российской Федерации;	\item \textit{природе разрыва} (какая часть обуславливается экономическим и культурным капиталом).
	\end{enumerate}
	
	\section{Обзор литературы}\label{sec:litreview}
	Социологические исследования российского фермерства установили, что фермеры представляют собой отдельную, но хрупкую социальную страту в сельской классовой структуре. Великий П.П. \cite{velikiy2025fermerstvo} показывает, что российские фермеры не являются однородной группой, а стратифицированы по размеру хозяйства: крупные фермеры наиболее прочно «укоренены» в агропромышленном комплексе, тогда как число глав малых хозяйств сократилось вдвое в период с 2006 по 2016 год . Ключевой тезис Великого состоит в том, что воспроизводство фермерской группы находится под угрозой из-за отсутствия приемников, желающих продолжать фермерство: дети фермеров, как правило, мигрируют в города . Эта внутренняя хрупкость свидетельствует о том, что фермеры, несмотря на их статус наиболее благополучного сегмента сельского населения, действуют в условиях структурной неопределенности, отличающей их от других профессиональных групп.
	
	Теория капитала и поля Бурдье предоставляет продуктивную рамку для интерпретации этого структурного положения. Бурдье различает экономический, культурный и социальный капитал и утверждает, что ценность любого капитала определяется специфическим полем, в котором действуют акторы. Эмпирические исследования отдачи от образования в России подтверждают эту структурную интерпретацию. Исследование А.В. Корицкого и С.Д. Капелюк \cite{koritsky2014urbanization} на данных РМЭЗ-ВШЭ показало, что в период с 2006 по 2011 год частная норма отдачи от образования для сельских работников была выше, чем для городских (в 2006 и 2009 гг. разница была статистически значима на уровне 5\%). Однако эта отдача была подвержена значительным колебаниям (от 3,7\% до 7,2\%), в то время как для городских работников она оставалась стабильной (4,2–4,7\%). Такая волатильность свидетельствует о структурной нестабильности сельскохозяйственного поля: отдача от культурного капитала фермеров не является устойчивой и зависит от внешних шоков (например, кризиса 2009 года), что соответствует бурдьезианскому тезису о несоответствии полей.
	
	\section{Методология}\label{sec:methods}
	\subsection{Данные}
	Для анализа различий культурного и экономического капитала был применен метод декомпозиции Блайндера-Оаксаки. Исследование опирается на данные 34-й волны Российского мониторинга экономического положения и здоровья населения НИУ ВШЭ (РМЭЗ-ВШЭ), проводимой в 2025 году. РМЭЗ-ВШЭ представляет собой ежегодное лонгитюдное обследование, репрезентативное на национальном уровне, которое традиционно используется для анализа доходов, занятости и социальной стратификации в России. 
	
	Помимо РМЭЗ-ВШЭ существуют микроданные сельскохозяйственной микропереписи 2021 года (СХМП-2021)\footnote{Росстат: оф. страница СХМП-2021. URL: https://rosstat.gov.ru/shmp/2021}. Микроданные описывают категории хозяйств: сельскохозяйственные организации, крестьянские (фермерские) хозяйства (КФХ), индивидуальные предприниматели (ИП), личные подсобные хозяйства (ЛПХ) граждан и некоммерческие объединения. Также данные включают в себя показатели общей земельной площади, используемых сельхозугодий, посевных площадей и поголовья скота. Однако данные находятся в закрытом формате и не доступны без дополнительных соглашений с Росстатом, что делает этот источник временно недоступным на момент проведения исследования.
	
	Для формирования аналитической выборки из общей базы были исключены респонденты с нулевыми или отсутствующими значениями доходов, а также лица, не имеющие весовых коэффициентов. Итоговая выборка составила 3 343 наблюдений, из которых 68 респондентов были идентифицированы как фермеры. Группирующая переменная «фермер» (\textit{is\_farmer}) была сконструирована на основе самоидентификации респондента (переменная \textit{J90} = 10), что соответствует подходу, примененному в работе Великого П. П. \cite{velikiy2024fermerstvo} для выделения фермерства как особой социально-профессиональной группы. Однако в отличии от подхода Великого П. П. в выборку были добавлены лица работающие в агрикультурной сфере, самозанятые и продающие сельскохозяйственную продукцию, но при этом не идентифицировавшие себя как фермеры в вопросе об основной работе \textit{J90}. Дополнительно из выборки были исключены домохозяйки и студенты, которые также занимаются продажей сельскохозяйственной продукции. В качестве зависимой переменной используется логарифм личного дохода за последние 30 дней (переменная \textit{J60}), предварительно винзорированного на 95-м процентиле для устранения влияния экстремальных выбросов.
	
	Контрольная группа была сконструирована из жителей городов, урбанизированных районов и деревень в качестве противовеса основному критерию отбора: наличия статуса фермера. Выборка из жителей урбанизированных районов может привести к  \textit{смещению вследствие пропущенных переменных} в силу большого количество пропущенных значений в переменной\textit{ i3 (тип населенного пункта).} При этом спилловер-эффект не распространяется на обладателей дач, торгующих сельскохозяйственной продукцией так как этот доход не является основным и существенным. Такие наблюдения были отсеяны из целевой выборки.
	
	Для оценки детерминант дохода используются уравнения минцеровского типа, оцениваемые отдельно для фермеров и контрольной группы. Спецификация модели включает возраст и его квадрат (для учета нелинейной отдачи от опыта), пол, наличие высшего образования, семейное положение и размер домохозяйства. Включение социально-демографических переменных способствует стабильности модели и в классификации Бурдье каждая переменная может объясняться согласно классификации в таблице ~\ref{tab:bourdieu_classification}. 
	
	
	\begin{table}[htbp]
		\centering
		\caption{Классификация переменных в рамках теории капитала Бурдье}
		\label{tab:bourdieu_classification}
		\begin{tabular}{l l}
			\toprule\toprule
			Переменная & Форма капитала \\
			\midrule
			Общий доход за 30 дней 
			& Экономический капитал \\
			\addlinespace
			
			Наличие высшего образования 
			& Институционализированный культурный капитал \\
			\addlinespace
			
			Возраст 
			& Воплощённый культурный капитал (опыт) \\
			\addlinespace
			
			Семейное положение
			& Социальный капитал (связующий) \\
			\addlinespace
			
			Размер семьи 
			& Социальный капитал / Экономический капитал \\
			\addlinespace
			
			Пол
			& часть Габитуса (как социально-демографическая характеристика) \\
			\addlinespace
			
			\midrule
			Необъясненная часть 
			& Несоответствие «поля» и разница невключенных факторов \\
			\bottomrule
		\end{tabular}
	\end{table}
	
	В таблице~\ref{tab:bourdieu_classification}  представлена интепретация переменных в рамках теории Бурдье. В частности стоит отметить, что \textit{пол} в теории Бурдье представляется как фундаментальное, воплощенное измерение габитуса, которое определяет, как индивиды воспринимают социальный мир, проживают его и ориентируются в нем. Бурдье различает эффект пола и гендера в социальном пространстве. Пол может восприниматься как часть символического насилия, когда неравенство в доходах или привилегиях может быть объяснена через символическое насилие (мягкая, невидимая форма принуждения, осуществляемая с молчаливого согласия). В исследовании пол интерпретируется как часть габитуса и социально-демографическая характеристика.
	
	Таблица ~\ref{tab:desc_stats_ttl} описывает статистические характеристики всей совокупности.
	\begin{table}[htbp]
		\centering
		\caption{Описательная статистика всей выборки (взвешенные оценки)}
		\label{tab:desc_stats_ttl}
		\begin{tabular}{l r r r r r r r}
			\toprule\toprule
			Переменная & \textit{N} & \textit{Взвешенная N} & \textit{Средняя} & \textit{Ср. Откл.} & \textit{Медиана} & \textit{Мин} & \textit{Макс} \\
			\midrule
			Общий доход за 30 дней      & 8810 & 8390.04 & 48184.84 & 32724.82 & 40000 & 200 & 200000 \\
			Возраст                     & 8810 & 8390.04 & 49.88    & 18.01    & 48.5  & 14  & 98.5   \\
			Пол (женский)               & 8810 & 8390.04 & 0.55     & 0.50     & 1     & 0   & 1      \\
			Размер семьи                & 8810 & 8390.04 & 3.47     & 2.04     & 3     & 1   & 15     \\
			Наличие высшего образования & 8810 & 8390.04 & 0.30     & 0.46     & 0     & 0   & 1      \\
			Семейное положение (женат)  & 8810 & 8390.04 & 0.58     & 0.49     & 1     & 0   & 1      \\
			Фермер                      & 8810 & 8390.04 & 0.01     & 0.09     & 0     & 0   & 1      \\
			\bottomrule
		\end{tabular}
	\end{table}
	
	Таблица~\ref{tab:desc_stats_farmers} описывает характеристики выборки фермеров из общей панели РМЭЗ-ВШЭ. 
	\begin{table}[htbp]
		\centering
		\caption{Описательная статистика всей выборки (взвешенные оценки)}
		\label{tab:desc_stats_farmers}
		\begin{tabular}{l r r r r r r r}
			\toprule\toprule
			Переменная & \textit{N} & \textit{Взвешенная N} & \textit{Средняя} & \textit{Ср. Откл.} & \textit{Медиана} & \textit{Мин} & \textit{Макс} \\
			\midrule
			Общий доход за 30 дней      & 68 & 76.09 & 55304.16 & 28776.75 & 49900 & 12300 & 173000 \\
			Возраст                     & 68 & 76.09 & 44.04    & 11.80    & 44.0  & 21.5  & 66.5   \\
			Пол (женский)               & 68 & 76.09 & 0.21     & 0.41     & 0     & 0     & 1      \\
			Размер семьи                & 68 & 76.09 & 4.63     & 1.65     & 4     & 1     & 9      \\
			Наличие высшего образования & 68 & 76.09 & 0.34     & 0.47     & 0     & 0     & 1      \\
			Семейное положение (женат)  & 68 & 76.09 & 0.65     & 0.48     & 1     & 0     & 1      \\
			Фермер                      & 68 & 76.09 & 1.00     & 0.00     & 1     & 1     & 1      \\
			\bottomrule
		\end{tabular}
	\end{table}
	
	Важно заметить, что выборка фермеров в основном представлена людьми, состоящими в отношениях. В среднем 65\% фермеров вне зависимости от пола состоит в браке. Выборка представлена также четвертью женщин. Средний размер семьи выборки составляет 4.6 человека. Высшее образование имеют только 34\%. Средний возраст фермера в выборке составляет 44 года. Приложение Б дает наглядное сравнение основных характеристик фермеров и нефермеров. 
	
	\subsection{Модель}
	Для оценки детерминант личного дохода используются уравнения минцеровского типа, оцениваемые отдельно для фермеров и контрольной группы. Базовая спецификация модели имеет следующий вид:
	
	\begin{equation}
		\ln(Y_i + 1) = \beta_0 + \beta_1 \text{Age}_i + \beta_2 \text{Age}_i^2 + \beta_3 \text{Female}_i + \beta_4 \text{Educ}_i + \beta_5 \text{Married}_i + \beta_6 \text{FamilySize}_i + \varepsilon_i,
		\label{eq:mincer}
	\end{equation}
	где:
	\begin{itemize}
	\item $\ln(Y_i + 1)$ - натуральный логарифм личного дохода респондента $i$ за последние 30 дней; 
	\item $\text{Age}_i$ и $\text{Age}_i^2$ - возраст и его квадрат, отражающие нелинейную отдачу от опыта; 
	\item $\text{Female}_i$ - дамми-переменная пола (1 = женщина); 
	\item $\text{Educ}_i$ - наличие высшего образования; 
	\item $\text{Married}_i$ - дамми-переменная семейного положения; 
	\item $\text{FamilySize}_i$ - размер домохозяйства.
	\end{itemize}
	
	Оценка уравнения \eqref{eq:mincer} производится методом взвешенных наименьших квадратов (WLS) и робастными стандартными ошибками (HC3), что позволяет корректировать гетероскедастичность и обеспечивать репрезентативность результатов. Следует отметить, что малый размер подвыборки фермеров ($N = 68$) накладывает ограничения на количество включаемых регрессоров. В частности, включение региональных фиксированных эффектов было признано нецелесообразным ввиду угрозы потери степеней свободы и нестабильности оценок. Поэтому модель ограничивается индивидуальными характеристиками, которые, согласно теории капитала и поля Бурдье, отражают объём культурного и социального капитала, доступного актору в конкретном поле.
	
	Для разложения разрыва в доходах между фермерами и контрольной группой на объяснённую и необъяснённую части применяется метод декомпозиции Блайндера-Оаксаки \cite{blinder1973wage, oaxaca1973male}. Пусть $\bar{Y}_F$ и $\bar{Y}_C$ обозначают средние натуральные логарифмы доходов фермеров и контрольной группы соответственно, а $\bar{X}_F$ и $\bar{X}_C$ - векторы средних значений объясняющих переменных. Тогда общий разрыв в доходах может быть представлен в следующем виде:
	\begin{equation}
		\bar{Y}_F - \bar{Y}_C
		= \underbrace{(\bar{X}_F-\bar{X}_C)'\beta_C}_{\text{эндаументы}}
		+ \underbrace{\bar{X}_C'(\beta_F-\beta_C)}_{\text{коэффициенты}}
		+ \underbrace{(\bar{X}_F-\bar{X}_C)'(\beta_F-\beta_C)}_{\text{взаимодействие}},
		\label{eq:oaxaca}
	\end{equation}
	где
	\begin{itemize}
		\item \(\beta_F\) и \(\beta_C\)~- векторы коэффициентов из
		регрессий для фермеров и контрольной группы соответственно;
		\item \(\bar{X}_F\) и \(\bar{X}_C\)~- векторы средних значений
		объясняющих переменных в соответствующих группах.
	\end{itemize}
	
	Эффект эндаументов в уравнении \eqref{eq:oaxaca} представляет собой
	\textit{объяснённую часть}. Она представлена долей разрыва,
	связанная с различиями в наблюдаемых характеристиках (образовании, возрасте, семейном положении, размере домохозяйства). Эндаументы оцениваются по отдачам контрольной групп	\(\beta_C\). В терминах теории капитала Бурдье эта часть отражает различия в объеме институционализированного и воплощённого культурного капитала, а также связующего социального капитала между группами.

	Эффект коэффициентов представляет собой \textit{необъясненную часть}. Необъясняемая часть может трактоваться как указание на «несоответствие поля»: одинаковый культурный капитал (например, диплом о высшем образовании)
	приносит разную отдачу в сельскохозяйственном и несвязанных с ним полях. Эффект отражает различия в отдаче от одинаковых характеристик при фиксированной композиции контрольной группы.  Однако необъяснённая часть аккумулирует не только различия в отдачах, но и влияние ненаблюдаемых характеристик (предпринимательские способности, доступ к земле, кредитам и субсидиям), региональную гетерогенность и ошибки измерения, поэтому она только \emph{совместима} с тезисом о несоответствии поля, но не доказывает его.
	
	Эффект взаимодействия возникает потому, что композиционные различия и различия в отдачах действуют одновременно и не аддитивно. В терминах бурдьезианской теории взаимодействие не имеет самостоятельной теоретической интерпретации, но представляет из себя статистический артефакт совместного действия композиции и отдач.
	
	\section{Результаты}\label{sec:results}
	
	В таблице~\ref{tab:ols_compare} представлены результаты оценивания уравнения минцеровского типа \eqref{eq:mincer} отдельно для всей выборки, для подвыборки фермеров и для контрольной группы. Для контрольной группы (столбец 3) все коэффициенты при объясняющих переменных статистически значимы на уровне 1\%. Возраст положительно связан с доходом ($\beta = 0.0603$, $p < 0.001$), а квадратичный член отрицателен ($\beta = -0.0006$, $p < 0.001$), что отражает классический профиль накопления и последующего обесценивания воплощённого культурного капитала. Образование демонстрирует сильную положительную отдачу. Наличие высшего образования увеличивает логарифм дохода на 0.3516 ($p < 0.001$). Женщины зарабатывают значительно меньше мужчин ($\beta = -0.2064$, $p < 0.001$). Семейное положение и размер домохозяйства также оказывают значимое влияние.
	
	\begin{table}[H]
		\caption{Сравнение результатов регрессионного анализа}
		\label{tab:ols_compare}
		\begin{center}
			\begin{tabular}{llll}
			\toprule\toprule
			& (1) total  & (2) farmers & (3) control  \\
			\midrule
			Интерцепт        & 9.0924***  & 9.9744***   & 9.1946***    \\
			& (0.0995)   & (0.7535)    & (0.0972)     \\
			Возраст              & 0.0657***  & 0.0103      & 0.0603***    \\
			& (0.0037)   & (0.0370)    & (0.0035)     \\
			$\mathrm{Возраст}^{2}$     & -0.0006*** & -0.0001     & -0.0006***   \\
			& (0.0000)   & (0.0004)    & (0.0000)     \\
			Пол (женский)       & -0.2067*** & 0.0654      & -0.2064***   \\
			& (0.0168)   & (0.1551)    & (0.0166)     \\
			Наличие в. образ. & 0.3512***  & 0.0883      & 0.3516***    \\
			& (0.0151)   & (0.1399)    & (0.0147)     \\
			Сем. положение  & 0.1335***  & 0.3586**    & 0.1152***    \\
			& (0.0164)   & (0.1516)    & (0.0158)     \\
			Размер семьи     & -0.0224*** & 0.0587      & -0.0215***   \\
			& (0.0051)   & (0.0464)    & (0.0049)     \\
			\midrule
			R-squared        & 0.1891     & 0.1733      & 0.1772       \\
			R-squared Adj.   & 0.1886     & 0.0919      & 0.1766       \\
			F-statistic      & 276.7144***   & 2.3887*      & 261.8089***     \\
			F-stat. p-value    & 0.000   & 0.0387      & 0.000     \\
			N                & 8810       & 68          & 8742         \\
			\bottomrule
			\end{tabular}
		\end{center}
	\end{table}
	
	Противоположная картина наблюдается в подвыборке фермеров (столбец 2). Модель статистически значима ($F = 2.3887$, $p = 0.0387$), однако большинство индивидуальных коэффициентов не достигают статистической значимости. В частности, образование не оказывает значимого влияния на доход фермеров ($\beta = 0.0883$, $p = 0.5303$), в отличие от контрольной группы, где отдача от образования высока и значима. Возраст и его квадрат также незначимы ($p = 0.7822$ и $p = 0.8491$ соответственно), что свидетельствует об отсутствии классического профиля накопления опыта у фермеров. Единственной значимой переменной остается семейное положение ($\beta = 0.3586$, $p = 0.0212$), что может отражать роль социального капитала (отражающихся в виде связующих сетей) в сельскохозяйственной деятельности. Эти результаты являются свидетельством в пользу бурдьезианского тезиса о несоответствии полей: культурный капитал фермеров не конвертируется в экономический доход в рамках сельскохозяйственного поля.
	
	В таблице~\ref{tab:oaxaca_results} представлены результаты трехчастной декомпозиции Блайндера-Оаксаки, оценённой методом объединённой регрессии (декомпозиция Оаксаки-Блайндера). Общий разрыв в логарифмах доходов, определённый как \(\Delta=\overline{\ln Y}_F-\overline{\ln Y}_C\), где \(F\)~- фермеры, \(C\)~- контрольная группа, составляет \(0.2776\) лог-пункта.
	Это означает, что средний логарифм дохода фермеров выше, чем у контрольной группы, примерно на \(0.2776\) лог-пункта, что соответствует примерно \(32\%\) более высокому геометрическому среднему доходу (\(\exp(0.2776)-1\approx 0.320\)).
	Это означает, что фермеры в среднем имеют более высокий логарифм дохода. 
	
	\begin{table}[htbp]
		\centering
		\caption{Результаты декомпозиции Блайндера-Оаксаки}
		\label{tab:oaxaca_results}
		\begin{tabular}{lrr}
			\toprule\toprule
			Компонент & Оценка & \% от разрыва \\
			\midrule
			Общий разрыв (\(\Delta \bar{Y}\)) & 0.2776 & 100.0\% \\
			\midrule
			Объяснённая часть (эндаументы) & 0.1578 & 56.8\% \\
			Необъяснённая часть (коэффициенты) & 0.1714 & 61.7\% \\
			Эффект взаимодействия & -0.0516 & -18.6\% \\
			\bottomrule
			\multicolumn{3}{l}{\footnotesize Примечание: классическая трехчастная
				декомпозиция Блайндера-Оаксаки.} \\
		\end{tabular}
	\end{table}
	
	Объясненная часть разрыва составляет \(0.1578\) лог-пункта, или \(56.8\%\) от общего разрыва. Поскольку общий разрыв определен как \(\Delta=\overline{\ln Y}_F-\overline{\ln Y}_C\) и положителен, положительное значение объясненной части означает, что наблюдаемые характеристики фермеров при эталонных отдачах связаны с более высоким	доходом, чем характеристики контрольной группы. Иными словами, по совокупности наблюдаемых признаков (возрастной структуры, образования, семейного положения, размера домохозяйства и пола) фермеры в среднем находятся в более благоприятном положении. 
	
	Однако доминирующим компонентом разрыва является необъясненная часть,
	которая составляет \(0.1714\) лог-пункта, или \(61.7\%\) от общего разрыва. При выбранном направлении разрыва (\(F-C\)) положительное значение этого компонента означает, что отдачи от наблюдаемых характеристик и неучтенные факторы в среднем благоприятствуют фермерам в большей степени, чем контрольной группе. Если бы фермеры обладали теми же наблюдаемыми характеристиками, что и
	контрольная группа, остаточная компонента всё равно смещала бы разрыв
	в пользу фермеров. В рамках теории капитала и поля Бурдье отдачи от одних и тех же
	характеристик различаются между полями: сельскохозяйственное поле вознаграждает культурный и социальный капитал иначе, чем другие поля.
	Отрицательный эффект взаимодействия (\(-0.0516\), или \(-18.6\%\)) указывает на то, что композиционный и коэффициентный компоненты частично компенсируют друг друга: более благоприятная наблюдаемая композиция фермеров сочетается с
	относительно менее выраженным преимуществом в отдачах, так что итоговый
	разрыв оказывается меньше суммы первых двух компонентов.
	
	\begin{table}[H]
		\caption{Результаты декомпозиции Оаксаки-Блайндера (трехкомпонентная)}
		\label{tab:oaxaca_blinder}
		\begin{center}
			\begin{tabular}{lrrrrrrrr}
				\toprule\toprule
				& \multicolumn{2}{c}{Эндаумент} & \multicolumn{2}{c}{Коэффициент} & \multicolumn{2}{c}{Взаимодействие} & \multicolumn{2}{c}{Всего} \\
				\cmidrule(lr){2-3}\cmidrule(lr){4-5}\cmidrule(lr){6-7}\cmidrule(lr){8-9}
				Переменная & Знач. & \% & Знач. & \% & Знач. & \% & Знач. & \% \\
				\midrule
				Возраст                  & -3.0014 & -1081.3 &  1.6763 &  603.9 &  0.2581 &   93.0 & -1.0671 & -384.4 \\
				$\mathrm{Возраст}^{2}$   &  3.9183 &  1411.7 & -2.1924 & -789.9 & -0.8662 & -312.1 &  0.8597 &  309.7 \\
				Размер семьи             &  0.2969 &   107.0 &  0.2499 &   90.0 & -0.0655 &  -23.6 &  0.4813 &  173.4 \\
				Сем. положение           & -0.0050 &    -1.8 &  0.2465 &   88.8 & -0.0364 &  -13.1 &  0.2051 &   73.9 \\
				Наличие в. образ.        &  0.0031 &     1.1 & -0.1091 &  -39.3 &  0.0116 &    4.2 & -0.0944 &  -34.0 \\
				Пол (женский)            & -0.0654 &   -23.6 &  0.0260 &    9.4 &  0.0366 &   13.2 & -0.0027 &   -1.0 \\
				\midrule
				Итого                    &  1.1465 &   413.1 & -0.1028 &  -37.0 & -0.6618 & -238.4 &  0.3819 &  137.6 \\
				\bottomrule
			\end{tabular}
		\end{center}
		
		\vspace{0.5em}
		\footnotesize
		Группа: is\_farmer (0 vs 1), направление: 1 $-$ 0 (фермеры $-$ контрольная группа).\\
		Средние исходы: группа 0 = 10.5160, группа 1 = 10.7936, разность = 0.2776.
	\end{table}
	
	Наиболее содержательным в свете теории Бурдье является вклад высшего образования. Эндаумент-эффект практически равен нулю (\(0.0031\)), что отражает близость долей лиц с высшим образованием в двух группах (\(34\%\) у фермеров против \(30\%\) в контрольной группе). Однако коэффициентный компонент отрицателен (\(-0.1091\)) и составляет около \(-39\%\) от общего разрыва: более низкая отдача от образования в сельскохозяйственном поле (\(0.0883\) у фермеров против \(0.3516\) в контроле) сокращает преимущество фермеров примерно на \(0.11\) лог-пункта. Такое положение отражает прямое эмпирическое соответствие бурдьезианскому тезису о несоответствии поля: институционализированный культурный капитал (диплом) приносит в сельскохозяйственном поле существенно меньшую отдачу, чем в других полях. Однако в силу размера выборки коэффициент статистически незначительный.
	
	Второй содержательно интерпретируемый компонент: семейное положение. Эндаумент-эффект вновь близок к нулю (\(-0.0050\)), поскольку доля состоящих в браке различается умеренно (\(65\%\) у фермеров против \(58\%\) в контроле). Напротив, коэффициентный компонент велик и положителен (\(0.2465\)), что соответствует \(+88.8\%\) от общего разрыва: отдача от брака у фермеров (\(0.3586\), \(p=0.018\)) значительно выше, чем в контрольной группе (\(0.1152\)), и этот дифференциал дает фермерам преимущество порядка \(0.25\) лог-пункта. В терминах бурдьезианского подхода это можно прочитать как свидетельство роли связующего социального капитала: семейные сети в сельскохозяйственном поле конвертируются в экономический доход эффективнее, чем в городских полях, где занятость в большей степени опосредована формальными институтами.
	
	Вклад размера семьи следует интерпретировать с осторожностью: эндаумент-эффект положителен (\(0.2969\)), что отражает больший средний размер домохозяйства фермеров (\(4.63\) против \(3.47\)). Качественно результат совместим с интерпретацией семьи как производственного и социального ресурса в фермерском хозяйстве, но количественно на него опираться не следует в силу отсутствия статистической значимости эффекта. Гендерный компонент мал по абсолютной величине (эндаумент \(-0.0654\), коэффициент \(0.0260\), взаимодействие \(0.0366\)). Такая ситуация отражает как низкую долю женщин среди фермеров (\(21\%\)), так и отсутствие значимого гендерного разрыва в фермерской подвыборке.
	
	\section{Ограничения}\label{sec:discussion}
	Настоящее исследование имеет ряд ограничений, которые необходимо учитывать при интерпретации полученных результатов. Наиболее существенное ограничение связано с малым размером подвыборки фермеров ($N = 68$). Хотя модель для фермеров является статистически значимой ($F = 2.3887$, $p = 0.0387$), большинство индивидуальных коэффициентов не достигают уровня статистической значимости, что ограничивает возможность содержательной интерпретации отдельных детерминант дохода. Малое число наблюдений также не позволило включить в модель региональные фиксированные эффекты, которые могли бы контролировать различия в стоимости жизни и структуре сельскохозяйственных рынков между субъектами РФ. В результате оценки могут быть смещены из-за ненаблюдаемой региональной гетерогенности, а необъясненная часть разрыва может частично отражать эти различия, а не только «несоответствие поля» в терминах бурдьезианского подхода.
	
	Также существует ограничение связанное с кросс-секционным характером данных, который не позволяет установить причинно-следственные связи. Декомпозиция Блайндера-Оаксаки описывает статистическую ассоциацию между характеристиками и доходами, однако точную причину, является ли низкая отдача от образования для фермеров следствием структурной дискриминации или результатом самоотбора в фермерскую деятельность, определить невозможно в рамках исследования. Кроме того, необъясненная часть разрыва может содержать влияние ненаблюдаемых переменных, таких как риск-предпочтения, предпринимательские способности, доступ к земельным ресурсам или институциональным субсидиям, которые не измеряются в РМЭЗ-ВШЭ. Наконец, использование самоидентификации (переменная \textit{J90} = 10) для определения фермеров может привести к недооценке группы, поскольку часть фермеров может идентифицировать себя как «самозанятые» или «предприниматели».
	
	\section{Заключение}\label{sec:conclusion}
	

	%они поддерживают более слабый тезис о том, что отдачи от одних и тех же характеристик различаются между группами.
	
	%Этот результат \emph{совместим} с бурдьезианским тезисом о различии отдач от одного и того же культурного капитала в разных полях, но не является его прямым подтверждением. Широкий доверительный интервал для фермеров включает и нулевое, и экономически значимое положительное значение; следовательно, вывод об отсутствии отдачи от образования в сельскохозяйственном поле при \(N=68\) сделать нельзя. Корректная формулировка: точечная оценка для фермеров существенно ниже, чем для контрольной группы, однако из-за малого размера подвыборки она статистически неотличима от нуля.
	
	%сельскохозяйственное поле вознаграждает культурный и социальный капитал иначе, чем другие поля, 
	
	Исследование позволило охарактеризовать фермеров как социально-профессиональную группу и оценить природу разрыва в доходах между фермерами и остальным населением. Декомпозиция Блайндера-Оаксаки показала, что средний логарифм дохода фермеров превышает аналогичный показатель контрольной группы на \(0.2776\) лог-пункта, что соответствует примерно \(32\%\) преимуществу по геометрическому среднему. При этом \(56.8\%\) разрыва объясняются различиями в наблюдаемых характеристиках (эндаументами), тогда как \(61.7\%\) приходятся на необъясненную часть, отражающую различия в отдачах от этих характеристик. При этом отрицательный эффект взаимодействия (\(-0.0516\); \(-18.6\%\)) частично компенсирует сумму первых двух компонентов. Покомпонентная декомпозиция (таблица~\ref{tab:oaxaca_blinder}) локализует этот коэффициентный канал: он формируется главным образом разнонаправленными вкладами отдачи от образования (отрицательный для фермеров: \(-0.1091\)) и отдачи от семейного положения (положительный: \(+0.2465\)).
	
	Полученные результаты совместимы с бурдьезианским тезисом о несоответствии полей. Отдача от одних и тех же характеристик систематически различаются между сельскохозяйственным полем и другими полями. В частности, точечная оценка отдачи от высшего образования для фермеров (\(\beta = 0.0883\)) существенно ниже, чем для контрольной группы (\(\beta = 0.3516\)), однако из-за малого размера подвыборки (\(N = 68\)) она статистически неотличима от нуля. Полученный результат не позволяет сделать категоричный вывод об отсутствии отдачи от культурного капитала в сельскохозяйственном поле, но указывает на необходимость его дальнейшей проверки на репрезентативных данных.
	
	Содержательно, экономическое положение фермеров определяется не столько дефицитом культурного или экономического капитала, сколько структурной организацией сельскохозяйственного рынка труда. Наблюдаемые преимущества фермеров в композиции характеристик сочетаются с менее выраженным преимуществом в отдачах. Меры социальной политики, направленные исключительно на наращивание образовательного потенциала сельского населения, могут оказаться недостаточными без изменения институциональных условий, определяющих конвертацию культурного и социального капитала в экономический доход.
	
	\begin{thebibliography}{7}
		
		\bibitem{bokontaeva1998farmers}
		Д.~К. Боконтаева.
		\newblock Фермеры как социально-профессиональная группа.
		\newblock \emph{СОЦИС}, 1998.
		
		\bibitem{velikiy2024fermerstvo}
		П.~П. Великий.
		\newblock Фермерство - (не)укоренившаяся социальная группа: постановка проблемы.
		\newblock \emph{Социологические исследования}, №5. С. 79-88, 2024.
		\newblock DOI: 10.31857/S0132162524050064.
		\newblock URL: https://www.socis.isras.ru/article/10060 
		
		\bibitem{velikiy2025fermerstvo}
		П.~П. Великий.
		\newblock Фермерство: жизненный мир и социальное самочувствие.
		\newblock \emph{Социологические исследования}, №2. С. 29-39, 2025.
		\newblock DOI: 10.7868/S3034601025020037.
		\newblock URL: https://www.socis.isras.ru/article/10397
		
		\bibitem{koritsky2014urbanization}
		А.~В. Корицкий, С.~Д. Капелюк.
		\newblock Влияние урбанизации на частную и социальную отдачу от образования в России.
		\newblock \emph{XIV Апрельская международная научная конференция по проблемам развития экономики и общества}, 2014.
				
		\bibitem{bourdieu1986forms}
		P.~Bourdieu.
		\newblock The Forms of Capital.
		\newblock In J.~G. Richardson, editor, \emph{Handbook of Theory and Research for the Sociology of Education}, Greenwood Press, New York, 1986.
		
		\bibitem{blinder1973wage}
		A.~S. Blinder.
		\newblock Wage Discrimination: Reduced Form and Structural Estimates.
		\newblock \emph{Journal of Human Resources}, 1973.
		
		\bibitem{oaxaca1973male}
		R.~Oaxaca.
		\newblock Male-Female Wage Differentials in Urban Labor Markets.
		\newblock \emph{International Economic Review}, 1973.
		
		
		
	\end{thebibliography}
	\newpage
	\appendix
	\section*{Приложение А. Взвешенная регрессионная модель фермеров}
		\begin{table}[H]
			\caption{Результат взвешенной линейной регрессии по выборке фермеров}
			\label{tab:ols_farmers}
			\begin{center}
				\begin{tabular}{llll}
					\hline
					Model:              & WLS                   & Adj. R-squared:     & 0.092     \\
					Dependent Variable: & np.log1p(ttl\_income) & AIC:                & 98.2563   \\
					Date:               & 2026-10-19 12:40      & BIC:                & 113.7928  \\
					No. Observations:   & 68                    & Log-Likelihood:     & -42.128   \\
					Df Model:           & 6                     & F-statistic:        & 2.389     \\
					Df Residuals:       & 61                    & Prob (F-statistic): & 0.0387    \\
					R-squared:          & 0.173                 & Scale:              & 0.24780   \\
					\hline
				\end{tabular}
			\end{center}
			
			\begin{center}
				\begin{tabular}{lrrrrrr}
					\hline
					&   Coef. & Std.Err. &       t & P$> |$t$|$ &  [0.025 &  0.975]  \\
					\hline
					Intercept        &  9.9744 &   0.7535 & 13.2382 &      0.0000*** &  8.4678 & 11.4811  \\
					age              &  0.0103 &   0.0370 &  0.2777 &      0.7822 & -0.0637 &  0.0842  \\
					$\mathrm{age}^{2}$      & -0.0001 &   0.0004 & -0.1911 &      0.8491 & -0.0010 &  0.0008  \\
					is\_female       &  0.0654 &   0.1551 &  0.4220 &      0.6745 & -0.2446 &  0.3755  \\
					education\_level &  0.0883 &   0.1399 &  0.6311 &      0.5303 & -0.1915 &  0.3681  \\
					marital\_status  &  0.3586 &   0.1516 &  2.3661 &      0.0212** &  0.0555 &  0.6617  \\
					family\_size     &  0.0587 &   0.0464 &  1.2660 &      0.2103 & -0.0340 &  0.1515  \\
					\hline
				\end{tabular}
			\end{center}
			
			\begin{center}
				\begin{tabular}{llll}
					\hline
					Omnibus:       & 0.094  & Durbin-Watson:    & 1.752  \\
					Prob(Omnibus): & 0.954  & Jarque-Bera (JB): & 0.121  \\
					Skew:          & -0.078 & Prob(JB):         & 0.941  \\
					Kurtosis:      & 2.865  & Condition No.:    & 35922  \\
					\hline
				\end{tabular}
			\end{center}
		\end{table}
			
		\section*{Приложение Б. Сравнение фермеров и нефермеров}
		\begin{figure}[H]
			\centering
			\includegraphics[width=0.8\textwidth]{farmers_vs_controls_updated.png}
			\caption{Сравнение фермеров и нефермеров}
			\label{fig:infiographics}
		\end{figure}
	
\end{document}